\documentclass[10pt,a4paper]{amsart}
\usepackage[margin=19mm]{geometry}
\usepackage{amsmath,amssymb,mathtools}
\usepackage[scr=rsfs,cal=euler]{mathalfa}
\usepackage{xfrac}
\usepackage[unicode,hidelinks,bookmarks=false]{hyperref}
\newcommand{\ie}{i.e.\ }
\newcommand{\cf}{cf.\ }
\newcommand{\resp}{respectively\ }
\newcommand{\Subsection}{\subsection}
\providecommand{\nobreakdash}{\nobreak\hbox{-}}
\setlength{\emergencystretch}{2em}
\multlinegap0pt
\usepackage{algorithmic}
\usepackage{textcomp}
\def\BibTeX{{\rm B\kern-.05em{\sc i\kern-.025em b}\kern-.08em
    T\kern-.1667em\lower.7ex\hbox{E}\kern-.125emX}}
\hypersetup{pdftex,colorlinks=true,allcolors=blue} %comment for original
\usepackage{verbatim}
\usepackage{array}
\usepackage{mathrsfs}
\usepackage{enumitem}
\usepackage{lmodern}
\newtheorem{theorem}{Theorem}%[section]
\newtheorem{lemma}{Lemma}
\newtheorem{proposition}{Proposition}
\newtheorem{corollary}{Corollary}
\newtheorem{definition}{Definition}
\newcommand{\R}{\mathbb{R}}
\newcommand{\E}{\mathbb{E}}
\newcommand{\Cov}{{\rm Cov}}
\newcommand{\Var}{{\rm Var}}
\newcommand{\I}{\mathbf{1}}
\newcommand{\intd}{{\rm d}}
\DeclareMathOperator*{\argmax}{arg\,max}
\DeclareMathOperator*{\argmin}{arg\,min}
\DeclareMathOperator*{\arginf}{arg\,inf}
\def\cA{{\mathcal A}}
\def\cE{{\mathcal E}}
\def\cG{{\mathcal G}}
\def\cI{{\mathcal I}}
\def\cL{{\mathcal L}}
\def\cM{{\mathcal M}}
\def\cP{{\mathcal P}}
\def\cS{{\mathcal S}}
\def\cV{{\mathcal V}}
\def\sA{{\mathsf A}}
\def\sL{{\mathsf L}}
\def\sM{{\mathsf M}}
\def\sN{{\mathsf N}}
\def\sS{{\mathsf S}}
\def\sU{{\mathsf U}}
\def\sV{{\mathsf V}}
\def\sW{{\mathsf W}}
\def\sX{{\mathsf X}}
\def\sY{{\mathsf Y}}
\def\sZ{{\mathsf Z}}
\def\bmZ{{\text{\boldmath$Z$}}}
\def\bmz{{\text{\boldmath$z$}}}
\def\rd{{\rm d}}
\def\PP{{\mathbb P}}
\def\QQ{{\mathbb Q}}
\def\N{{\mathbb N}}
\def\Alg{{\mathfrak A}}
\def\deq{\triangleq}
\def\wh#1{{\widehat{#1}}}
\def\eps{\varepsilon}
\begin{document}
\title[Minimum Distortion Quantization with Specified Output Distri]{Minimum Distortion Quantization with Specified Output Distribution}
\author{Aolin Xu}
\date{}
\maketitle
\noindent\emph{Source and licence.} Minimum Distortion Quantization with Specified Output Distribution. Version arXiv:2606.10458v1 (CC BY 4.0). This material is distributed under the Creative Commons Attribution 4.0 International licence, http://creativecommons.org/licenses/by/4.0/, as recorded in the arXiv OAI licence metadata for this exact version and shown on the abstract page https://arxiv.org/abs/2606.10458v1. Full rights notice: licenses/arxiv-2606.10458-CC-BY-4.0.txt.\par\smallskip
\noindent\textit{Editorial scope.} Counted unit: the original \texttt{theorem} environment \texttt{th:main} at source lines 437--448 together with its complete proof at lines 449--484; the delivered document keeps source lines 160--485 so that every referenced label and every citation resolves inside it. Native editable source is the paper's own concatenated TeX; the AMS wrapper is editorial formatting only. No publisher class, upstream script or unrelated artwork is redistributed.
\par\smallskip\noindent\textit{Reference note.} This article ships no compiled bibliography (biblatex); the entries delivered here are composed from the author's own BibTeX (.bib) fields, with no field invented and no entry dropped.
\begin{align}
    \min_{\substack{f:\,\R \rightarrow \{1,\ldots,k\} \\ \text{s.t. } f(W)\sim P_X}} \text{mmse}(W|f(W)) \label{eq:opt1}
\end{align}
where 
\begin{align}
\text{mmse}(W|f(W)) \deq \min_{g:\,\{1,\ldots,k\}\rightarrow\R}\E[(W-g(f(W)))^2]    
\end{align}
for any $f$.
Given the number of values $k$ that can be taken by the quantization output, restricting the range of the quantizer to be 
$\{1,\ldots,k\}$ is of no loss of generality, as there is a one-to-one mapping between $\{1,\ldots,k\}$ and any finite set of cardinality $k$, and such a mapping does not change the MMSE of estimating $W$ from the quantization output \cite[Lemma~1]{MER2022}.


\section{Derivation of optimal quantizer}
\subsection{Problem conversion}
For any quantizer $f:\R\rightarrow\{1,\ldots,k\}$ that maps $W$ to $X=f(W)$, using basic properties of MMSE and conditional expectation, it can be shown that
% (a proof is given in Appendix~\ref{sec:pf_mmse_decomp})
\begin{align}
    \text{mmse}(W|X) = \E[W^2]- \E[\E[W|X]^2] .\label{eq:mmse_decomp}
\end{align}
As $\E[W^2]$ is determined by $P_W$, whenever $\E[W^2] < \infty$, the optimization problem in \eqref{eq:opt1} is equivalent to 
\begin{align}
    \max_{\substack{f:\,\R \rightarrow \{1,\ldots,k\} \\ \text{s.t. } f(W)\sim P_X}} \E[\E[W|f(W)]^2] . \label{eq:max_E[m^2]}
\end{align}
Under the assumption that $W$ is a continuous random variable, $F_W$ is invertible and we can define
\begin{align}
    Q_W(u) \deq F_W^{-1}(u) , \quad u\in[0,1] .
\end{align}
With a random variable $U$ uniformly distributed over $[0,1]$, the tuple $(Q_W(U), f(Q_W(U))$ has the same joint distribution as $(W, f(W))$ for any $f:\R\rightarrow\{1,\ldots,k\}$, hence
\begin{align}
 \E[\E[W|f(W)]^2] = \E[\E[Q_W(U)|f(Q_W(U))]^2] .
\end{align}
For notational simplicity, we represent $P_X$ as
\begin{align}
    p_i \deq P_X(i) , \quad i=1,\ldots,k .
\end{align}
We have the following lemma.
\begin{lemma}\label{lm:convert_f_S}
The maximum of \eqref{eq:max_E[m^2]} satisfies
\begin{align}
& \max_{\substack{f:\,\R \rightarrow \{1,\ldots,k\} 
\\ 
\textnormal{s.t. } f(W)\sim P_X}} \E[\E[W|f(W)]^2] = 
\nonumber \\
& 
\max_{\substack{\sS_1,\ldots,\sS_k\subset[0,1] 
\\
\textnormal{s.t. } 
\sS_i\cap \sS_j = \emptyset, \, \cup_{i=1}^k \sS_i = [0,1] \\
P_U(\sS_i)=p_i } }
\sum_{i=1}^k p_i \, \left(\frac{1}{p_i}\int_{\sS_i}Q_W(u) \rd u \right)^2 . \label{eq:max_S}
\end{align}
Moreover, an $f^*$ satisfying the constraint and achieving the maximum on the LHS of \eqref{eq:max_S} can be obtained as
\begin{align}
    f^*(w) = i \quad\textnormal{if $w\in Q_W(\sS^*_i)$} , \quad w\in\R \label{eq:f*_def}
\end{align}
where $\sS^*_1,\ldots,\sS^*_k$ are subsets of $[0,1]$ satisfying the constraints and achieving the maximum on the RHS of \eqref{eq:max_S}, and $Q_W(\sS^*_i) \deq \{Q_W(u): u\in \sS^*_i\}$.

\begin{proof}
For any $f:\R\rightarrow\{1,\ldots,k\}$ satisfying $f(W) \sim P_X$,
\begin{align}
& \E[\E[W|f(W)]^2] \nonumber \\
=& \E[\E[Q_W(U)|f(Q_W(U))]^2]  \\
=& 
\sum_{i=1}^k \PP[f(Q_W(U))=i] \, \E[Q_W(U)|f(Q_W(U))=i]^2 \\
=& 
\sum_{i=1}^k p_i \, \left(\frac{1}{p_i}\int_{\sS_i}Q_W(u) \rd u \right)^2 \label{eq:E_E[W|f(W)]^2}
\end{align}
where
\begin{align}
    \sS_i \deq \{u \in [0,1]: f(Q_W(u))=i\} . \label{eq:S_def}
\end{align}
As $f$ induces an ordered partition $(\sW_1,\ldots,\sW_k)$ of $\R$, which may also be called a labeled partition,
with 
\begin{align}
\sW_i \deq \{w\in\R: f(w)=i\}
\end{align} 
and $(F_W, Q_W)$ is a pair of bijections between $\R$ and $[0,1]$, the $k$-tuple of subsets $(\sS_1,\ldots,\sS_k)$ of $[0,1]$ defined in \eqref{eq:S_def} constitute an ordered partition of $[0,1]$, satisfying
$\sS_i = F_W(\sW_i)$, $\sS_i\cap \sS_j = \emptyset$, $\cup_{i=1}^k \sS_i = [0,1]$ and $P_U(\sS_i)=p_i$.
This proves that
\begin{align}
& \max_{\substack{f:\,\R \rightarrow \{1,\ldots,k\} 
\\ 
\textnormal{s.t. } f(W)\sim P_X}} \E[\E[W|f(W)]^2] 
\le \nonumber \\
& 
\max_{\substack{\sS_1,\ldots,\sS_k\subset[0,1] 
\\
\textnormal{s.t. } \sS_i\cap \sS_j = \emptyset, \, \cup_{i=1}^k \sS_i = [0,1], \\
P_U(\sS_i)=p_i}}
\sum_{i=1}^k p_i \, \left(\frac{1}{p_i}\int_{\sS_i}Q_W(u) \rd u \right)^2 . \label{eq:max_f_ineq}
\end{align}
Next is to show that the equality in \eqref{eq:max_f_ineq} can be achieved by $f^*$ defined in \eqref{eq:f*_def}.
To this end, let $\sS^*_1,\ldots,\sS^*_k$ be subsets of $[0,1]$ satisfying the constraints and achieving the maximum on the RHS of \eqref{eq:max_f_ineq}.
For $f^*$ defined as
\begin{align}
    f^*(w) = i \quad\textnormal{if $w\in Q_W(\sS^*_i)$} , \quad w\in\R , \nonumber
\end{align}
it maps $\R$ to $\{1,\ldots,k\}$ and $\PP[f(W)=i]=\PP[W\in Q_W(\sS^*_i)] = \PP[Q_W(U)\in Q_W(\sS^*_i)] = \PP[U\in\sS^*_i] = p_i$, hence $f^*$ satisfies the constraint on the LHS of \eqref{eq:max_f_ineq}.
Moreover, for this $f^*$,
\begin{align}
% $
   \{u: f^*(Q_W(u))=i\} 
   = \{u: Q_W(u) \in Q_W(\sS^*_i)\} 
  =\sS^*_i .
% $ .
\end{align}
Together with \eqref{eq:E_E[W|f(W)]^2}, it follows that $f^*$ achieves the RHS of \eqref{eq:max_f_ineq}. This proves the lemma.
\end{proof}

\end{lemma}

Lemma~\ref{lm:convert_f_S} converts the quantizer optimization problem in \eqref{eq:max_E[m^2]}
to an ordered partition optimization problem of the unit interval on the RHS of \eqref{eq:max_S}.
Its proof also shows a one-to-one correspondence between any admissible quantizer $f$ satisfying the constraint on the LHS of \eqref{eq:max_S} and an admissible ordered partition $(\sS_1,\ldots,\sS_k)$ satisfying the constraints on the RHS of \eqref{eq:max_S} through \eqref{eq:S_def} and $\eqref{eq:f*_def}$.

\subsection{Solution to converted problem}
The converted optimization problem through Lemma~\ref{lm:convert_f_S} can be restated as
\begin{align}
\max_{\substack{\sS_1,\ldots,\sS_k\subset[0,1] 
\\
\textnormal{s.t. } \sS_i\cap \sS_j = \emptyset, \, \cup_{i=1}^k \sS_i = [0,1], \\
P_U(\sS_i)=p_i}}
\sum_{i=1}^k p_i \, m_i^2 \label{eq:converted_max_Em^2}
\end{align}
where
\begin{align}
    m_i &\deq \frac{1}{p_i}\int_{\sS_i}Q_W(u) \rd u 
    = \E[Q_W(U)|U\in\sS_i].
\end{align}
For the $f:\R\rightarrow\{1,\ldots,k\}$ induced by an admissible ordered partition $(\sS_1,\ldots,\sS_k)$ of $[0,1]$ through \eqref{eq:f*_def}, 
\begin{align}
    m_i = \E[W|f(W)=i] .
\end{align}
The set of all admissible ordered partitions 
\begin{align}
\{(\sS_1,\ldots,\sS_k): 
& \, \sS_i\subset [0,1],  \sS_i\cap \sS_j = \emptyset,  \nonumber \\
&  \cup_{i=1}^k \sS_i = [0,1], P_U(\sS_i)=p_i\} \label{eq:def_set_ord_part}
\end{align}
% $\{(\sS_1,\ldots,\sS_k): \sS_i\cap \sS_j = \emptyset, \, \cup_{i=1}^k \sS_i = [0,1], P_U(\sS_i)=p_i\}$ 
can be divided into $k!$ categories according to the order of the values of $m_1, \ldots, m_k \in \R$ induced by each $(\sS_1,\ldots,\sS_k)$.
Specifically, each category can be represented by a permutation $\sigma$ of $\{1,\ldots,k\}$, such that 
\begin{align}
    m_{\sigma(1)} \le \ldots \le m_{\sigma(k)} \label{eq:order_m}
\end{align}
for each $(\sS_1,\ldots,\sS_k)$ in that category.
We have the following lemma.
\begin{lemma}\label{lm:major}
Among the category of the admissible ordered partitions in \eqref{eq:def_set_ord_part} resulting in $m_{\sigma(1)} \le \ldots \le m_{\sigma(k)}$, the one defined by contiguous intervals
\begin{align}
    \sS^*_{\sigma(i)} \deq [q_{i-1}, q_i)  \quad i=1,\ldots,k \label{eq:def_S*_q}
\end{align}
with
\begin{align}
    q_i \deq q_{i-1} + p_{\sigma(i)}
\end{align}
and $q_0 \deq 0$ satisfies
\begin{align}
    \sum_{i=1}^j p_{\sigma(i)} m^*_{\sigma(i)}
    \le
    \sum_{i=1}^j p_{\sigma(i)} m_{\sigma(i)} \quad j=1,\ldots, k-1 \label{eq:major_j}
\end{align}
and 
\begin{align}
    \sum_{i=1}^k p_{\sigma(i)} m^*_{\sigma(i)}
    =
    \sum_{i=1}^k p_{\sigma(i)} m_{\sigma(i)} \label{eq:major_k}
\end{align}
for all other ordered partitions in this category.
\end{lemma}
\begin{proof}
First, by the construction of $\sS^*_{\sigma(i)}$, the order in \eqref{eq:order_m} is met, as $Q_W(u)$ is increasing.
For $j=1,\ldots,k$, we have
\begin{align}
\sum_{i=1}^j p_{\sigma(i)} m^*_{\sigma(i)}
&= \sum_{i=1}^j \int_{\sS^*_{\sigma(i)}} Q_W(u) \rd u \\
&= \int_{\cup_{i=1}^j \sS^*_{\sigma(i)}} Q_W(u) \rd u
\end{align}
and for any other $(\sS_1,\ldots,\sS_k)$ in this category
\begin{align}
\sum_{i=1}^j p_{\sigma(i)} m_{\sigma(i)}
&= \sum_{i=1}^j \int_{\sS_{\sigma(i)}} Q_W(u) \rd u \\
&= \int_{\cup_{i=1}^j \sS_{\sigma(i)}} Q_W(u) \rd u .
\end{align}
Let $\sA_j \deq \cup_{i=1}^j \sS^*_{\sigma(i)}$ and ${\mathsf B}_j \deq \cup_{i=1}^j \sS_{\sigma(i)}$. Then $P_U(\sA_j) = P_U(\mathsf B_j)$, hence
\begin{align}
    P_U(\sA_j\setminus\mathsf B_j) = P_U(\mathsf B_j \setminus \sA_j) . \label{eq:B-A=A-B}
\end{align}
It follows that
\begin{align}
& \sum_{i=1}^j p_{\sigma(i)} m_{\sigma(i)} 
-
\sum_{i=1}^j p_{\sigma(i)} m^*_{\sigma(i)} \nonumber \\
=& 
\int_{\mathsf B_j} Q_W(u) \rd u 
-
\int_{\sA_j} Q_W(u) \rd u \\
=& \int_{\mathsf B_j\setminus\sA_j} Q_W(u) \rd u 
-
\int_{\sA_j\setminus\mathsf B_j} Q_W(u) \rd u \\
\ge& P_U(\mathsf B_j \setminus \sA_j) \Big( \inf_{u\in \mathsf B_j \setminus \sA_j} Q_W(u) - \sup_{u\in \mathsf A_j \setminus \mathsf B_j} Q_W(u)\Big) \label{eq:major_B-A} \\
\ge & 0  \label{eq:major_B>A}
\end{align}
where \eqref{eq:major_B-A} follows from \eqref{eq:B-A=A-B} and the fact that $Q_W(u)$ is increasing in $u$, and \eqref{eq:major_B>A} follows from the fact that $Q_W(u)$ is increasing in $u$ and 
\begin{align}
    \sup( \sA_j\setminus\mathsf B_j )
\le 
    \inf( \mathsf B_j\setminus\sA_j )
\end{align}
by the construction of $\sS^*_{\sigma(i)}$ for $i=1,\ldots,j$ in \eqref{eq:def_S*_q}.
This proves \eqref{eq:major_j}.
The proof of \eqref{eq:major_k} is straightforward as
\begin{align}
    \sum_{i=1}^k p_{\sigma(i)} m^*_{\sigma(i)}
    =
    \sum_{i=1}^k p_{\sigma(i)} m_{\sigma(i)} 
    = \int_{[0,1]} Q_W(u) \rd u  = \E[W]. \nonumber
\end{align}
\end{proof}
Lemma~\ref{lm:major} shows that $(m^*_{\sigma(1)}, \ldots , m^*_{\sigma(k)})$ induced by the ordered partition $(\sS^*_1,\ldots,\sS^*_k)$ defined in \eqref{eq:def_S*_q} weighted-majorizes $( m_{\sigma(1)}, \ldots , m_{\sigma(k)})$ induced by any other ordered partition $(\sS_1,\ldots,\sS_k)$ in the same category.
The following lemma can be viewed as a weighted majorization inequality, a.k.a.\ weighted Karamata inequality \cite{Fuchs_weighted, Majorization_inequalities_IT}. 
We include a proof for completeness.
\begin{lemma}\label{lm:wkata}
For $m^*_1 \le \ldots \le m^*_k \in \R$, $m_{1} \le \ldots \le m_{k} \in \R$, and $p_1,\ldots, p_k \in \R$, if
\begin{align}
    \sum_{i=1}^j p_{i} m^*_{i}
    \le
    \sum_{i=1}^j p_{i} m_{i} \quad j=1,\ldots, k-1 \label{eq:wkata_j}
\end{align}
and 
\begin{align}
    \sum_{i=1}^k p_{i} m^*_{i}
    =
    \sum_{i=1}^k p_{i} m_{i} \label{eq:wkata_k}
\end{align}
then for any convex function $\varphi: \R \rightarrow \R$,
\begin{align}
    \sum_{i=1}^k p_{i} \varphi(m^*_{i})
    \ge
    \sum_{i=1}^k p_{i} \varphi( m_{i} ) .
\end{align}
\end{lemma}
\begin{proof}
For $j=1,\ldots,k$, define
\begin{align}
a_j \deq \sum_{i=1}^j p_i m^*_i, \quad
b_j \deq  \sum_{i=1}^j p_i m_i
\end{align}
and $a_0 = b_0 \deq 0$.
Also, define
\begin{align}
c_i \deq \frac{\varphi(m_i) - \varphi(m^*_i)}{m_i - m^*_i} .
\end{align}
The convexity of $\varphi$ implies that
\begin{align}
    c_i \le c_{i+1} . \label{eq:wkata_c}
\end{align}
We have
\begin{align}
    & \quad \sum_{i=1}^k p_{i} \varphi(m^*_{i})
    -
    \sum_{i=1}^k p_{i} \varphi( m_{i} ) \nonumber \\
    % &= \sum_{i=1}^k c_i(p_i m_i - p_i m^*_i) \\
    &= \sum_{i=1}^k c_i (a_i - a_{i-1} - b_i + b_{i-1}) \\
    &= \sum_{i=1}^k c_i (a_i - b_i)
    - \sum_{i=1}^k c_i (a_{i-1} - b_{i-1}) \\
    &= c_k(a_k-b_k) + \sum_{i=1}^{k-1}(c_i - c_{i+1})(a_i - b_i) \label{eq:wkata_1}\\ %- c_1 (a_0 - b_0) \\
    &\ge 0 . \label{eq:wkata_2}
\end{align}
where \eqref{eq:wkata_1} is due to $a_0=b_0=0$, and \eqref{eq:wkata_2} follows from the facts that $a_k-b_k=0$ by \eqref{eq:wkata_k}, $a_i - b_i \le 0$ by \eqref{eq:wkata_j}, and $c_i - c_{i+1} \le 0$ by \eqref{eq:wkata_c}.
\end{proof}
The conventional majorization inequality is a special case of Lemma~\ref{lm:wkata} by taking $p_1=\ldots=p_k=1$ and reversing the indices $(1,\ldots,k)$. 
Together with Lemma~\ref{lm:major}, by taking $\varphi(m)=m^2$, Lemma~\ref{lm:wkata} proves the in-category optimality of the ordered partition $(\sS^*_1,\ldots,\sS^*_k)$ defined in \eqref{eq:def_S*_q}, as stated in the following lemma.
\begin{lemma}\label{lm:in-catagory}
Among the category of the admissible ordered partitions in \eqref{eq:def_set_ord_part} resulting in $m_{\sigma(1)} \le \ldots \le m_{\sigma(k)}$, the one defined by contiguous intervals in \eqref{eq:def_S*_q} achieves the maximum of $\sum_{i=1}^k p_{i} m^{2}_{i}$.
\end{lemma}
With Lemma~\ref{lm:convert_f_S} and Lemma~\ref{lm:in-catagory} in hand, we can state the main result of this work.

\begin{theorem}\label{th:main}
The optimal quantizer that achieves the minimum in \eqref{eq:opt1} can take the form
\begin{align}
X = \sigma^*\big(F_{\sigma^{*-1}(X)}^{-1}(F_W(W))\big) \label{eq:thm_f}
\end{align}
where $\sigma^*$ is the permutation of $\{1,\ldots,k\}$ that maximizes
\begin{align}
    \sum_{i=1}^k \frac{1}{p_i} \Big(\int_{\sS_i}Q_W(u) \rd u \Big)^2 \label{eq:thm_Em^2}
\end{align}
with $\sS_{\sigma(i)} = [q_{i-1}, q_i)$, $q_i = q_{i-1} + p_{\sigma(i)}$, $i=1,\ldots,k$,
and $q_0 = 0$ among all permutations $\sigma$.
\end{theorem}

\begin{proof}
The analysis developed so far shows that the maximization on the RHS of \eqref{eq:max_S} can be achieved by first dividing the set of admissible ordered partitions of $[0,1]$ into $k!$ categories according to the order of $m_1, \ldots, m_k$, and then optimizing within each category.
It follows from Lemma~\ref{lm:in-catagory} that in a category represented by the permutation $\sigma$, the optimal ordered partition takes the form
$\sS_{\sigma(i)} = [q_{i-1}, q_i)$ with $q_i = q_{i-1} + p_{\sigma(i)}$, $i=1,\ldots,k$,
and $q_0 = 0$. 
The globally optimal ordered partition is then determined by the permutation $\sigma^*$ that maximizes \eqref{eq:thm_Em^2}.
It then follows from Lemma~\ref{lm:convert_f_S}
and \eqref{eq:max_E[m^2]} that the quantizer $f^*$ determined by $\sigma^*$ achieves the minimum $\text{mmse}(W|f^*(W))$ while satisfying $f^*(W)\sim P_X$.

It remains to show that this quantizer takes the particular form in \eqref{eq:thm_f}.
To see this, recall that the quantizer determined by $\sigma^*$ is given by \eqref{eq:f*_def} as
\begin{align}
    f^*(w) = \sigma^*(i) \quad\textnormal{if $w\in Q_W(\sS_{\sigma^*(i)})$} , \quad w\in\R . \label{eq:thm_pf_f*}
\end{align}
We have
\begin{align}
w\in Q_W(\sS_{\sigma^*(i)}) 
\Leftrightarrow F_W(w) \in 
\sS_{\sigma^*(i)} ,
\end{align}
$F_W(W)$ has the same distribution as $U$, and the $k$-tuple 
\begin{align}
& \quad 
(P_U(\sS_{\sigma^*(1)}), \ldots, P_U(\sS_{\sigma^*(k)})) 
 \nonumber \\ 
&= (P_{X}(\sigma^*(1)), \ldots, P_{X}(\sigma^*(k))) \\
&= (P_{\sigma^{*-1}(X)}(1), \ldots, P_{\sigma^{*-1}(X)}(k))
\end{align}
is the probability mass function of the random variable $\sigma^{*-1}(X)$, whose generalized inverse CDF maps the contiguous subintervals
$(\sS_{\sigma^*(1)}, \ldots, \sS_{\sigma^*(k)})$ of $[0,1]$ to the indices $(1,\ldots,k)$.
With these observations, we know the function in \eqref{eq:thm_pf_f*} is equivalent to
\begin{align}
f^*(w) = \sigma^*\big(F_{\sigma^{*-1}(X)}^{-1}(F_W(w))\big), \quad w\in\R .
\end{align}
This completes the proof of the theorem.
\end{proof}


\begin{thebibliography}{99}
\bibitem[]{Fuchs_weighted} L. Fuchs. A new proof of an inequality of {H}ardy-{L}ittlewood-{P}olya. Matematisk Tidsskrift B (1947) 53-54
\bibitem[]{MER2022} Xu, Aolin; Raginsky, Maxim. Minimum Excess Risk in {B}ayesian Learning. IEEE Transactions on Information Theory 68 (2022) 7935-7955
\bibitem[]{Majorization_inequalities_IT} MA Khan; SI Bradanovic; N Latif; D Pecaric; J Pecaric. Majorization Inequality and Information Theory. (2019)
\end{thebibliography}
\end{document}
